\section{Threats to Validity}
\label{sec:threats}

We organize the threats along the four standard validity
categories~\cite{cook1979quasi,wohlin2012experimentation}.
The supporting measurements and mechanics for each are in the online supplement; every limitation is
stated here.

\paragraph{Construct validity of the measures}
Our dependent variables are HTTP-level throughput (req/s) and tail latency (p99) at the Express
endpoint, not query-execution time at the driver boundary.
This is deliberate, capturing the full request round trip, though overhead attributed to the
\emph{access layer} may partly reflect Express routing and serialization.
We hold those near-constant: every adapter satisfies a documented \emph{adapter contract} and funnels
rows through shared canonical constructors whose standalone cost is too small to explain the
multi-fold gap, and a fixed-payload endpoint puts the shared framework-and-serialization floor at
least three times above the fastest deep-fetch cell, so no cell is framework-bound.
The specification-derived oracle and the differential checks establish finite conformance and mutual
equivalence (Section~\ref{sec:methodology}); the differences arise in how each layer reaches an
identical response, not in what is returned.
The reported single-row-insert finding uses each engine's \emph{default} vendor-standard durability,
not one we relaxed, so those insert spreads (Supplement Table~S52) describe the tested default
deployment rather than a tuning choice; the relaxed-durability sensitivity is a compound
intervention and does not decompose its cause (Section~\ref{sec:discussion}, Supplement Table~S3).

\paragraph{Construct validity of the treatment}
RQ1's treatment is each layer's \emph{policy-selected documented} implementation and eager-loading
strategy: an operational rule not validated against surveyed, typical, or performance-tuned usage.
It is an intentionally artificial, predeclared selection policy, not a behavioral practitioner model.
The frozen pages, conflict rule, and assignments make the decision auditable, but the single author
both defined and applied it, and two blind selection reviews are prepared and still pending.
Where the policy differs from a performance-tuned configuration, RQ1's spread mixes the library's
execution machinery with that strategy choice, made material by the one-to-four statement range on
the deep fetch (Supplement Table~S2).
Three sensitivity checks qualify it, and none of them decomposes the spread: the preserved
source-page provenance is provenance, not evidence of frequency or idiomaticity; the
\emph{common-SQL raw-path sensitivity contrast} is a compound contrast, not a bound on how much
comes from the loading strategy versus the library machinery; and the loading-strategy check shows
the selected path is the \emph{faster} strategy for the join-selected ORMs, with Objection's slower
selected select-in path on MySQL disclosed rather than corrected.
RQ1 therefore answers what this policy yields, not a library-only overhead.

\paragraph{Internal validity}
Each adapter follows the registered policy and passes specification-derived and differential checks,
but all were implemented by one author.
A semantically correct yet needlessly allocating, converting, or serializing adapter would pass; no
independent human implementation review is claimed.
A result-blind author self-audit found two avoidable timed-path operations before campaign
acceptance, both removed without changing the declared treatment or SQL, after which the admission
chain passed again and measurement restarted from repetition~1
(Section~\ref{sec:discussion}); because the implementer performed it, this is author review.

Both are now detected mechanically instead, which is why R7 requires evidence rather than a reviewer:
a static check reports values acquired and never read, and a runtime gate rejects any cell writing to
\texttt{stderr} or \texttt{stdout} during a measured run.
The runtime gate postdates the reported campaign and did not gate it, so absence of output there
rests on author observation.
Neither covers waste that is silent \emph{and} invisible in the source, so the residual narrows
without closing: a correct, policy-conformant adapter uniformly slower than a differently written one
would pass every check.
We report it as open, not discharged.
The same caution applies more widely: four admission elements postdate the accepted campaigns
(Supplement Table~S48).
Because they replay a deterministic seed and inspect archived state, they can invalidate archived
results, and did so for the superseded pilot; but they did not gate the timed runs as those runs
executed, so a defect observable only while a run was in progress would have escaped them.

Per-cell warm-up phases exceeding the slowest layer's cold-start stabilization absorb JIT, pool-fill,
and plan-cache effects~\cite{mytkowicz2009producing}, and two further controls address non-layer
confounds: one correlated-subquery aggregation plan and the corrected-state campaign's verified reset
before each write cell.
The historical MySQL allocator defect, the corrected reset, and the seeded identifier design are
described in Section~\ref{sec:methodology}; their consequence here is that the affected historical
cells are not used as clean evidence, all five patterns having been remeasured in one campaign.
A forward-versus-reversed write-cell-order check found no detectable order effect, so order does not
confound layer identity with run position.
A residual risk is coordinated omission, which can under-report the tail at
saturation~\cite{fruth2022tail,tene2015latency}; the corrected constant-arrival sweeps show small
tails below measured capacity and sharp growth near it (Supplement Tables~S6 and S17).

\paragraph{External validity}
All measurements come from a single schema, dataset size, and 16-vCPU virtualized host, a pool fixed
at ten, and specific engine and library versions (PostgreSQL~18.4, MySQL~9.7.1; latest compatible
stable versions as of the 15 July 2026 freeze, Supplement Table~S11).
The fixed pool does not appear to drive the ranking: the pool-size frontier preserves the ordering
from 2 to 50 connections and a mixed read/write workload preserves the read ordering, both
exploratory and on four layers, though at a pool of one Prisma overtakes Knex on PostgreSQL
(Supplement Tables~S7, S25, S26).
Two choices are disclosed rather than resolved: Sequelize on its stable 6.x line, and Prisma's MySQL
path on the MariaDB adapter, an implementation$\times$engine asymmetry.
The memory-resident working set is a hot-cache regime that makes access-layer overhead more visible;
as a workload becomes I/O-bound the spreads may compress toward $1\times$.

Same-host checks do not bound an independent substrate: a post-restart rerun of six representative
cells moved absolute throughput materially against the campaign prevailing when it ran, which is the
superseded pilot rather than the accepted campaign.
The corrected-state 25-run campaign repeats all five patterns with a new randomized order and
environment fingerprint on the same host.
Cross-campaign agreement or disagreement is reported directly; it reduces dependence on one overnight
ordering but does not address host, hypervisor-neighbor, CPU-steal, or frequency transfer.
No empirical ordering is intended to travel beyond the tested campaigns, and close cross-backend rank
reversals are exploratory unless their direction repeats.
Every promoted RQ2 reversal involves the single treatment whose MySQL driver differs, and a
leave-one-out repetition of the promotion procedure without it promotes none (Supplement Table~S49);
the promoted-reversal evidence is therefore conditional on that driver substitution, which the design
cannot separate from the treatment or from the DBMS.
The protocol evidence is also bounded: its external audit has one rater, so general independent
applicability remains to be tested.

\paragraph{Conclusion validity}
Our analysis (Section~\ref{sec:methodology}) reports medians with the coefficient of variation and
within-campaign bootstrap 95\% intervals (repeatability, not population-general) and, because the
design is a randomized block, compares adjacently ranked layers with \emph{paired} tests on their
per-replicate throughput ratios.
Two bounds reinforce the ordering: the widest spreads dwarf the deep-fetch run-to-run variability,
while insert variability is large enough on both engines that no fine insert ordering is claimed
(Supplement Tables~S1, S9), and the ordering holds across the concurrency sweep at $\geq32$
connections (Supplement Figure~S2).
Rather than assert a specific statistical power, absent a prespecified effect-size model, we rest the
conclusions on effect sizes and interval estimates, foremost the \emph{predefined} native-driver
contrasts (Supplement Table~S41), keeping the ranking-selected adjacent-pair permutation tests only
as a descriptive post-hoc check.
Application processes do not cross adapter, engine, or replicate blocks, although read endpoints
within a block share one process~\cite{kalibera2013rigorous}.
All blocks share one host; primary intervals remain within-campaign, while the corrected RQ2 subset
has a second same-host campaign.
The secondary experiments covering only a layer subset or few replicates (Supplement Table~S32) are
read as \emph{sensitivity evidence} on the tested subset, not as claims about all eleven
implementations.
The common-SQL contrast and the matched-utilization tails rest on few replicates, so both are
reported with their observed spread and neither carries an interval estimate (Supplement
Tables~S14, S21--S22).
The matched-utilization tails carry a second, unpropagated uncertainty: each offered rate is a
fraction of an \emph{estimated} capacity, whose spread is evaluated separately (Supplement
Table~S45) rather than carried into the reported p99.
They are secondary evidence for the operating-point distinction, not a precise tail estimate.
